\chapter{λ-calculus / type theory to Lean dictionary}\label{sec:root:lambda-calculus-dictionary:1}

A direct dictionary connecting the formal λ-calculus/type-theory notation used in a few ``Mathematical reading'' boxes throughout this book back to the Lean syntax it corresponds to. This page is a lookup table, not something to read start to finish --- \hyperref[sec:01-basics:04-terminology:1]{Chapter 1 §4} and \hyperref[sec:01-basics:05-pi-sigma-and-coc:1]{§5}, \hyperref[sec:03-propositions-and-proofs:02-logic-recap:1]{Chapter 3 §2}, and \hyperref[sec:05-rigor-check:03-typing-rules-and-safety:1]{Chapter 5 §3} are where each row is actually built up and explained.

\begin{tabular}{lll}
\toprule
λ-calculus / type theory & Lean syntax & Where in this book \\
\midrule
\bottomrule
variable \(x\) & \texttt{x} & Chapter 1 \\
abstraction \(\lambda x.\, t\) & \texttt{fun x => t} & Chapter 1 \\
application \(t_1\, t_2\) & \texttt{t1 t2} & Chapter 1 \\
α-conversion & (silent, automatic) & --- \\
β-reduction \((\lambda x.t)s \to t[x:=s]\) & \texttt{rfl}, \texttt{\#eval} reduction & Chapter 1 §4, Chapter 5 §4 (definitional equality) \\
type judgment \(\Gamma \vdash t : \tau\) & \texttt{\#check t} & Chapter 1 \\
simple function type \(\tau_1 \to \tau_2\) & \texttt{τ1 → τ2} & Chapter 1, Chapter 5 §3 \\
universe hierarchy \(\mathtt{Type}\,i\) & \texttt{Type}, \texttt{Type 1}, \ldots{} & Chapter 5 §2--§3 \\
Π-type \(\prod_{x:A} B(x)\) & \texttt{(x : A) → B x}, \texttt{∀ x : A, B x} & Chapter 1 §3, §5, Chapter 3 \\
Σ-type \(\sum_{x:A} B(x)\) & \texttt{structure} (extractable); \texttt{∃ x, P x} is a \emph{restricted} cousin (no witness-extraction) rather than literally Σ & Chapter 2, Chapter 1 §5 \\
proof-irrelevant universe & \texttt{Prop} & Chapter 3, Chapter 1 §5 \\
inductive data (beyond bare CoC) & \texttt{inductive}, \texttt{structure} & Chapters 1, 2, 11 \\
Curry--Howard: propositions-as-types & \texttt{Prop}/proof terms & Chapter 3 \\
progress + preservation & Lean's kernel type-checker & Chapter 5 §3 \\
Church numerals (encoded data) & \texttt{Nat} (but built via \texttt{inductive}, not encoded) & Chapter 1 §4, Chapter 13 (aside) \\
Church booleans / if-then-else via application & \texttt{Bool}, \texttt{if}/\texttt{match} (but native, not encoded) & Chapter 13 (aside) \\
\end{tabular}


Two related facts, explained where their own dictionary rows live rather than repeated here: tactics (\hyperref[sec:01-basics:04-terminology:1]{Chapter 1 §4}) add nothing to the underlying calculus --- they are a user interface for building the same terms step by step --- and elaboration/unification (\hyperref[sec:01-basics:04-terminology:1]{Chapter 1 §4}) is type inference for this calculus, a well-understood algorithm, not magic.
